\documentclass[11pt,a4paper]{article}
\usepackage[margin=2.5cm]{geometry}
\usepackage{amsmath,amssymb}
\usepackage{booktabs}
\usepackage[colorlinks=true,linkcolor=blue]{hyperref}
\newcommand{\wn}{\ensuremath{\mathrm{cm^{-1}}}}
\title{\vspace{-2cm}Electronic structure calculations: draft main-text section}
\date{}
\begin{document}
\maketitle
\vspace{-1.5cm}

\section{Computational results and their relation to the observed dynamics}

Electronic structure calculations were performed to identify the states reached
at the two excitation wavelengths and to determine which nuclear coordinates
they drive. All calculations used the TPSSh functional with a mixed
def2-TZVP(Fe,O)/def2-SVP(C,H) basis and SMD aqueous solvation, matching the
solvent of the experiment; the ground state was optimised without symmetry
constraints and confirmed as a true minimum. Excited states were obtained from
TDA-TD-DFT with 80 roots, and the quartet manifold from spin-flip TD-DFT with
140 roots. Full details, tables and the complete mode-by-mode analysis are given
in the Supporting Information.

\subsection{Both excitation windows access ligand-to-metal charge transfer}

The calculated spectrum places a state at 2.92~eV (425~nm) and a second at
4.52~eV (277~nm), matching the two experimental excitation windows. Both are
dominated by $\beta \to \beta$ excitations into the vacant Fe(III)
$\beta$-$d$ manifold and both retain sextet spin
($\langle S^2 \rangle = 8.80$ and 8.96 against 8.75 for a pure sextet), so both
are ligand-to-metal charge transfer in character.

The two are nevertheless distinguished by the nuclear coordinate they drive.
Projecting each excited-state gradient, evaluated at the ground-state geometry,
onto the ground-state normal modes gives the displacement each mode undergoes
on vertical excitation. For the lower state the largest activity is the
conjugated C=C/C=O ring stretch at 1560~\wn{}; for the higher state it is the
Fe--O chelate deformation at 435~\wn{}, with a Huang--Rhys factor of 1.94
against 0.80. The higher-energy excitation distorts the metal--ligand framework
two to three times more strongly, and deposits nearly twice the total
reorganisation energy (0.59 against 0.32~eV). The computed Fe--O mode at
435~\wn{} corresponds to the band observed near 449~\wn{}.

\subsection{Why the same three modes are populated at both wavelengths}

The reorganisation energy is concentrated in modes that the measurement does
not monitor. Summed over the $\nu_8$, $\nu_9$ and $\nu_{10}$ windows, the
Huang--Rhys factors amount to $S \approx 0.01$--0.04, and these bands together
carry only 2--6\% of the total resonance-Raman activity. The dominant
Franck--Condon activity lies instead at 150--450~\wn{} and at 1560~\wn{}.

The three observed bands are therefore not the primary distortion coordinates
of either excited state. They are reporters of the vibrational bath, populated
by whatever relaxation cascade is running above them. This accounts for the
observation that $\nu_8$, $\nu_9$ and $\nu_{10}$ appear with the same prompt
rise at both excitation wavelengths despite the two states differing by
1.6~eV in energy and by a factor of two in reorganisation energy: the bands
report the arrival of vibrational energy rather than the identity of the state
that delivered it.

\subsection{The 941~\wn{} transient band and excited-state symmetry breaking}

A charge-transfer excitation in a tris-chelate must lift the degeneracy of the
doubly degenerate ligand vibrations. In $D_3$ the ligand-framework modes
transform as $A_1$, $A_2$ and $E$, and the three equivalent chelate rings give
tight manifolds: the calculated $\nu_8$ region spans only 3.4~\wn{}, with an
$A$--$E$ separation of 1.3~\wn{}. A charge-transfer hole must be accommodated
by three equivalent rings, and wherever vibronic coupling exceeds that small
intrinsic splitting the delocalised arrangement is unstable and the hole
localises on one ring, as valence localises in a mixed-valence compound. The
distortion that achieves this destroys the equivalence on which the $E$
degeneracy rests.

The calculations confirm this on two independent grounds. Structurally, the six
Fe--O distances span 0.014~\AA{} in the ground state but 0.098 and 0.241~\AA{}
in the two excited states, the lower state showing two short and four long
distances, the signature of a single ring singled out. Vibrationally, all
twelve near-degenerate ground-state manifolds between 200 and 1700~\wn{} widen
in both excited states without exception, by median factors of 13 and 15 and by
no less than three. The $\nu_8$ manifold widens from 3.4 to 57~\wn{}, by
factors of 16.7 and 16.6 --- the same at both wavelengths to within 1~\wn{},
confirming that the effect follows from charge-transfer character rather than
from the properties of either state.

This bears directly on the transient band at 941~\wn{}, which appears at both
excitation wavelengths and which lies 63~\wn{} above the ground-state $\nu_8$
position. The calculations do not reproduce a rigid 63~\wn{} stiffening of
$\nu_8$: the individual descendants of that manifold shift by between $-43$
and $+24$~\wn{}. What they predict instead is that the manifold ceases to be a
single narrow feature and spreads over some 57~\wn{}, so that the transient
$\nu_8$ region should be substantially broader than the ground-state doublet at
878 and 896~\wn{} and may resolve into additional components. Since the
calculation already underestimates the ground-state splitting by a factor of
five (Supporting Information), the true excited-state spread is likely larger
still. On this reading the 941~\wn{} feature is best understood as a resolved
component of a split $\nu_8$ manifold rather than as a uniformly stiffened
mode, which also accounts for its appearance at both wavelengths.

\subsection{Why vibrational cooling is faster following 400~nm excitation}

The measured cooling times differ by a factor of four to five between the two
excitation wavelengths. This ordering does not follow from the properties of
either state considered alone; if anything those properties suggest the
opposite, since the higher state places more than twice the reorganisation
energy into the low-frequency skeletal modes that couple most efficiently to
the solvent. The difference instead reflects how many electronic states lie
below each excitation.

\begin{table}[htbp]
\centering
\caption{Electronic states below each excitation, from 80 sextet roots and 140
quartet roots computed at the ground-state geometry.}
\begin{tabular}{l r r}
\toprule
 & 2.92~eV (425~nm) & 4.52~eV (277~nm) \\
\midrule
Sextet states below              & 14 & 60 \\
Quartet states below             & 17 & 59 \\
\textbf{Total}                   & \textbf{31} & \textbf{119} \\
Nearest quartet / \wn{}          & 250 & 121 \\
Reorganisation energy / eV       & 0.32 & 0.59 \\
\bottomrule
\end{tabular}
\end{table}

The higher excitation has four times as many states beneath it, and the
measured cooling times differ by four to five. We propose that the measured
$\tau_2$ reflects the duration of the electronic cascade rather than the
intrinsic rate of energy transfer to the solvent. Each internal-conversion or
intersystem-crossing step converts an electronic gap into vibrational quanta,
so while the cascade is still descending it continuously resupplies vibrational
energy that the solvent is simultaneously removing. Following 400~nm excitation
2.9~eV is released through some 31 states and the cascade is short, so the
measured decay approaches the true cooling time. Following 270~nm excitation
4.5~eV is released through some 119 states, the cascade outlasts the cooling,
and the observed decay is set by the cascade.

The quartet manifold supplies a further asymmetry. The lower state lies in a
sparse region of the sextet manifold, 1049~\wn{} from the nearest same-spin
state, but only 250~\wn{} from a quartet --- comparable to $k_{\mathrm{B}}T$ and
small against the Fe(III) spin--orbit coupling constant. Same-spin internal
conversion is slow there and intersystem crossing should dominate, and because
the quartet manifold extends below the lowest sextet, a single spin-flip step
transfers the system onto a ladder that reaches further down than the sextet
manifold does. The higher state couples to the quartets more strongly still,
121~\wn{} to the nearest, but gains nothing by it: at 4.52~eV it retains 59
quartets beneath it. The shortcut is only useful to a state already near the
bottom of the manifold, and this is consistent with the faster sub-picosecond
decay of the 941~\wn{} feature observed following 400~nm excitation.

This explanation is a statistical one, resting on computed densities of states
rather than on calculated rates, and we present it as such. A quantitative
treatment would require spin--orbit couplings between the relevant sextet and
quartet pairs, from which intersystem-crossing rates could be evaluated using
the reorganisation energies reported here.

\end{document}
